% chapters/ch02_slutsky_equation_and_demand_response.tex

\section{Slutsky Equation and Demand Response}

Lecture 2 extends the revealed-preference framework from Chapter 1 and then turns to comparative statics of consumer demand. The chapter has two main aims. First, it strengthens WARP into a multi-observation consistency condition, GARP, which uses the transitivity of preferences. Second, it studies how Marshallian demand responds to changes in prices and expenditure, culminating in the Slutsky decomposition of a price effect into substitution and income effects.

\subsection{Generalized Revealed Preference}

Recall from Chapter 1 that if a bundle $x$ is chosen at prices $p$, then for any bundle $x'$:
\[
x R x' \iff p \cdot x \ge p \cdot x',
\qquad
x P x' \iff p \cdot x > p \cdot x'.
\]
If choices come from maximization of an increasing preference relation $\succeq$, then
\[
x R x' \implies x \succeq x',
\qquad
x P x' \implies x \succ x'.
\]

Because preferences are transitive, revealed preference should also propagate through chains of observations.

\begin{definition}{Revealed and strictly revealed preference through chains}{ch2-rstar}
Let $R$ and $P$ denote direct and strict direct revealed preference.

\begin{itemize}
    \item $x$ is \emph{revealed preferred} to $x'$, written $x R^* x'$, if there exist bundles
    \[
    x^1,\dots,x^N
    \]
    such that
    \[
    x R x^1 R x^2 R \cdots R x^N R x'.
    \]
    \item $x$ is \emph{strictly revealed preferred} to $x'$, written $x P^* x'$, if such a chain exists and at least one link is strict, i.e. at least one $R$ can be replaced by $P$.
\end{itemize}
\end{definition}

\begin{theorem}{Behavioral meaning of $R^*$ and $P^*$}{ch2-rstar-meaning}
Suppose observed choices maximize an increasing preference relation $\succeq$. Then
\[
x R^* x' \implies x \succeq x',
\qquad
x P^* x' \implies x \succ x'.
\]
\end{theorem}

\begin{proof}
If
\[
x R x^1 R x^2 R \cdots R x^N R x',
\]
then each link implies a weak preference comparison under utility maximization:
\[
x \succeq x^1,\quad x^1 \succeq x^2,\quad \dots,\quad x^N \succeq x'.
\]
By transitivity, $x \succeq x'$.

If at least one link is strict, then at least one strict preference comparison appears in the chain. Combining weak and strict comparisons by transitivity yields $x \succ x'$.
\end{proof}

\begin{definition}{Generalized Axiom of Revealed Preference (GARP)}{ch2-garp}
A dataset satisfies \emph{GARP} if
\[
x R^* x' \implies \text{not }(x' P x).
\]
Equivalently: once $x$ is revealed preferred to $x'$ through a chain of affordable choices, it cannot be that $x'$ is later chosen when $x$ was strictly cheaper.
\end{definition}

GARP is stronger than WARP because it checks for contradictions generated not only by one-step comparisons, but also by multi-step revealed-preference chains.

\begin{theorem}{Afriat's theorem, as used in this course}{ch2-afriat}
For a finite dataset of consumer choices:
\begin{enumerate}
    \item If the data are generated by maximization of an increasing utility function, then the data satisfy GARP.
    \item If the data satisfy GARP, then there exists an increasing utility function that rationalizes the data.
\end{enumerate}
Hence GARP is the correct revealed-preference test for rationalizability by an increasing utility function.
\end{theorem}

\begin{examtip}
In two-good settings, Lecture 2 notes that WARP and GARP are equivalent. In three or more goods, GARP is strictly stronger because cycles can arise through multiple observations even when no strict two-observation cycle exists.
\end{examtip}

\begin{examplebox}
\textbf{WARP can hold while GARP fails in three goods.}

Consider three observations:
\[
x^1=(1,0,0),\qquad x^2=(0,1,0),\qquad x^3=(0,0,1),
\]
at prices
\[
p^1=(2,4,1),\qquad p^2=(1,2,4),\qquad p^3=(4,1,2).
\]
Since each chosen bundle is a unit vector,
\[
p^1 \cdot x^1=2,\qquad p^2 \cdot x^2=2,\qquad p^3 \cdot x^3=2.
\]
Now compute cross-costs:
\[
p^1 \cdot x^3 = 1 \le 2 \implies x^1 R x^3,
\]
\[
p^3 \cdot x^2 = 1 \le 2 \implies x^3 R x^2,
\]
\[
p^2 \cdot x^1 = 1 \le 2 \implies x^2 R x^1.
\]
So
\[
x^1 R x^3 R x^2 R x^1,
\]
which gives a revealed-preference cycle. Moreover,
\[
p^2 \cdot x^1 = 1 < 2 = p^2 \cdot x^2,
\]
so in fact
\[
x^2 P x^1.
\]
Thus GARP is violated because $x^1 R^* x^1$ through a chain that ends with a strict reversal.

However, there is no strict \emph{two-observation} cycle, so WARP is not violated. This shows that GARP is strictly stronger than WARP in dimensions $L \ge 3$.
\end{examplebox}

\begin{commonmistake}
Do not test GARP by looking only for pairs $x R y$ and $y P x$. That checks WARP. GARP requires checking whether a contradiction appears \emph{after a chain} of revealed-preference comparisons.
\end{commonmistake}

\subsection{Budgets, Utility Maximization, and Demand}

The second part of the lecture returns to standard consumer theory. The two core objects are the utility function and the budget set.

\begin{definition}{Indifference curve and demand function}{ch2-demand-def}
Let $U : \mathbb{R}^L_+ \to \mathbb{R}$ be a utility function.

\begin{itemize}
    \item An \emph{indifference curve} is a set of bundles with the same utility level:
    \[
    \{x \in \mathbb{R}^L_+ : U(x)=\bar u\}.
    \]
    \item A \emph{demand function} is a map
    \[
    x : \mathbb{R}^L_{++} \times \mathbb{R}_{++} \to \mathbb{R}^L_+
    \]
    such that $x(p,m)$ is affordable and utility-maximizing at prices $p$ and expenditure $m$:
    \[
    p \cdot x(p,m) \le m,
    \qquad
    U(x(p,m)) \ge U(\tilde x)\ \ \forall \tilde x \in B(p,m).
    \]
\end{itemize}
\end{definition}

The demand function gives the consumer's optimal bundle as a function of prices and expenditure. Its $\ell$-th component, $x_\ell(p,m)$, is demand for good $\ell$.

\begin{theorem}{Exhausting the budget under increasing utility}{ch2-budget-exhaustion}
If $U$ is increasing and $x(p,m)$ maximizes $U$ over $B(p,m)$, then
\[
p \cdot x(p,m)=m.
\]
\end{theorem}

\begin{proof}
If $p \cdot x(p,m) < m$, then some strictly positive amount of some good can be added while remaining in the budget set because each price is strictly positive. The resulting bundle is strictly larger componentwise, so increasing utility implies it yields higher utility, contradicting optimality.
\end{proof}

\begin{examtip}
When the question gives a Marshallian demand function $x(p,m)$, always verify two things first:
\begin{enumerate}
    \item affordability: $p \cdot x(p,m) \le m$;
    \item if utility is increasing, budget exhaustion: $p \cdot x(p,m)=m$.
\end{enumerate}
These are the fastest internal consistency checks.
\end{examtip}

\subsection{Demand Response to Price Changes}

Let $x(p,m)=(x_1(p,m),x_2(p,m))$ denote Marshallian demand in the two-good case. The lecture studies what happens when the price of good 1 increases from $p_1$ to $p_1' > p_1$, holding $p_2$ and $m$ fixed.

\begin{definition}{Gross substitutes, gross complements, ordinary goods, and Giffen goods}{ch2-gross-def}
Suppose the price of good 1 rises from $p_1$ to $p_1'$.

\begin{itemize}
    \item If
    \[
    x_2(p_1',p_2,m) > x_2(p_1,p_2,m),
    \]
    then goods 1 and 2 are \emph{gross substitutes}.
    \item If
    \[
    x_2(p_1',p_2,m) < x_2(p_1,p_2,m),
    \]
    then goods 1 and 2 are \emph{gross complements}.
    \item If
    \[
    x_1(p_1',p_2,m) < x_1(p_1,p_2,m),
    \]
    then good 1 is \emph{ordinary}.
    \item If
    \[
    x_1(p_1',p_2,m) > x_1(p_1,p_2,m),
    \]
    then good 1 is \emph{Giffen}.
\end{itemize}
\end{definition}

These are purely Marshallian concepts: they describe the \emph{total} demand response to a price change before any decomposition into substitution and income effects.

\begin{examplebox}
\textbf{An explicit demand system: ordinary own-price effect and positive cross-price effect.}

Suppose
\[
U(x_1,x_2)=\sqrt{x_1}+\sqrt{x_2},
\]
which generates the Marshallian demand
\[
x(p_1,p_2,m)
=
\frac{m}{p_1+p_2}
\left(
\frac{p_2}{p_1},
\frac{p_1}{p_2}
\right).
\]
Hence
\[
x_1(p_1,p_2,m)=\frac{m p_2}{p_1(p_1+p_2)},
\qquad
x_2(p_1,p_2,m)=\frac{m p_1}{p_2(p_1+p_2)}.
\]

For good 2, holding $p_2$ and $m$ fixed,
\[
\frac{\partial x_2}{\partial p_1}
=
\frac{m}{p_2}\cdot \frac{p_2}{(p_1+p_2)^2}
=
\frac{m}{(p_1+p_2)^2}
>0.
\]
So when the price of good 1 rises, demand for good 2 rises: the goods are gross substitutes.

For good 1,
\[
\frac{\partial x_1}{\partial p_1}
=
-\frac{m p_2(2p_1+p_2)}{p_1^2(p_1+p_2)^2}
<0,
\]
so good 1 is ordinary, not Giffen.
\end{examplebox}

\begin{policynote}
The lecture's ``meat and potatoes'' example shows why Giffen behavior is theoretically possible. When a staple good becomes more expensive, the loss of real purchasing power may be so severe that the consumer buys \emph{more} of the staple and cuts back on the more expensive alternative. This is unusual, but it clarifies that the total Marshallian price effect need not always be negative.
\end{policynote}

\subsection{Demand Response to Expenditure Changes}

Now hold prices fixed and increase expenditure from $m$ to $m' > m$.

\begin{definition}{Normal and inferior goods}{ch2-normal-inferior}
A good $\ell$ is:
\begin{itemize}
    \item \emph{normal} if
    \[
    x_\ell(p,m') > x_\ell(p,m)
    \quad \text{whenever } m' > m;
    \]
    \item \emph{inferior} if
    \[
    x_\ell(p,m') < x_\ell(p,m)
    \quad \text{for some increase in expenditure } m' > m.
    \]
\end{itemize}
\end{definition}

Normality and inferiority concern income or expenditure effects, not substitution effects. A good can be ordinary or Giffen in price-response language, and normal or inferior in expenditure-response language.

\begin{examplebox}
\textbf{Normal goods in the square-root demand system.}

Continuing the example
\[
x_1(p_1,p_2,m)=\frac{m p_2}{p_1(p_1+p_2)},
\qquad
x_2(p_1,p_2,m)=\frac{m p_1}{p_2(p_1+p_2)},
\]
we compute
\[
\frac{\partial x_1}{\partial m}=\frac{p_2}{p_1(p_1+p_2)}>0,
\qquad
\frac{\partial x_2}{\partial m}=\frac{p_1}{p_2(p_1+p_2)}>0.
\]
Therefore both goods are normal.
\end{examplebox}

\begin{commonmistake}
Do not infer that a good is normal merely because its demand curve slopes downward. Downward-sloping Marshallian demand describes a price effect; normality describes the sign of the expenditure effect.
\end{commonmistake}

\subsection{Compensated Price Changes and the Slutsky Identity}

The central technical result of Lecture 2 is the decomposition of a price change into substitution and income effects.

Suppose prices change from
\[
p=(p_1,p_2)
\qquad \text{to} \qquad
p'=(p_1',p_2),
\quad p_1' > p_1,
\]
while original expenditure is $m$.

\begin{definition}{Slutsky compensated expenditure}{ch2-ms}
Let $x(p,m)$ be the original Marshallian demand. The \emph{Slutsky compensated expenditure} is
\[
m^s = p' \cdot x(p,m).
\]
This is the amount of money needed at the new prices to just afford the original bundle.
\end{definition}

So the consumer is ``compensated'' enough to remain able to purchase the original bundle after the price increase. This isolates a pure substitution comparison between the old chosen bundle and the new chosen bundle on a budget that passes through the old choice.

\begin{theorem}{The Slutsky substitution effect is nonpositive for the good whose price rises}{ch2-sub-neg}
If the price of good 1 rises from $p_1$ to $p_1' > p_1$, then
\[
x_1(p',m^s) - x_1(p,m) \le 0.
\]
\end{theorem}

\begin{proof}
By construction, the old bundle $x(p,m)$ is affordable at the new prices with compensated expenditure $m^s$ because
\[
p' \cdot x(p,m) = m^s.
\]
Hence $x(p',m^s)$ is chosen from a budget set that contains $x(p,m)$.

Suppose, for contradiction, that
\[
x_1(p',m^s) > x_1(p,m).
\]
Since the price of good 2 is unchanged, compare the old prices:
\begin{align*}
p \cdot x(p',m^s)
&= p_1 x_1(p',m^s) + p_2 x_2(p',m^s) \\
&< p_1' x_1(p',m^s) + p_2 x_2(p',m^s) \\
&= p' \cdot x(p',m^s) \\
&\le m^s \\
&= p' \cdot x(p,m) \\
&= p_1' x_1(p,m) + p_2 x_2(p,m).
\end{align*}
Because $x_1(p',m^s) > x_1(p,m)$ and $p_1' > p_1$, the only way this can hold is if the new choice is cheaper at old prices than the old choice, implying
\[
x(p,m) P x(p',m^s).
\]
But at the new compensated prices, the new choice is chosen while the old one is affordable, so
\[
x(p',m^s) R x(p,m).
\]
This violates WARP, and therefore also GARP, for utility-maximizing behavior. Hence the supposition was false, so
\[
x_1(p',m^s) \le x_1(p,m).
\]
\end{proof}

The lecture packages this into three objects.

\begin{keyformula}[title={Key Formula: Slutsky Decomposition}]
Let
\[
\Delta x_1 = x_1(p',m) - x_1(p,m)
\]
denote the total Marshallian change in demand for good 1 after the price increase. Define
\[
\Delta x_1^s = x_1(p',m^s) - x_1(p,m)
\]
as the \emph{Slutsky substitution effect}, and
\[
\Delta x_1^n = x_1(p',m) - x_1(p',m^s)
\]
as the \emph{Slutsky income effect}. Then
\[
\Delta x_1 = \Delta x_1^s + \Delta x_1^n.
\]
This is the \emph{Slutsky identity}.
\end{keyformula}

The identity is purely algebraic: add and subtract $x_1(p',m^s)$. Its power comes from the fact that the substitution effect has a definite sign, while the income effect depends on whether the good is normal or inferior.

\begin{theorem}{Law of demand for normal goods}{ch2-law-demand}
Suppose good 1 is normal. If its price rises from $p_1$ to $p_1' > p_1$, holding $p_2$ and $m$ fixed, then
\[
x_1(p_1',p_2,m) - x_1(p_1,p_2,m) \le 0.
\]
\end{theorem}

\begin{proof}
First,
\[
m^s = p' \cdot x(p,m)
= p_1' x_1(p,m) + p_2 x_2(p,m)
\ge
p_1 x_1(p,m) + p_2 x_2(p,m)
= m.
\]
So the compensated expenditure is at least as large as the original expenditure.

If good 1 is normal, then at the new prices $p'$ a higher expenditure raises demand for good 1:
\[
m^s \ge m
\implies
x_1(p',m^s) \ge x_1(p',m),
\]
which is equivalent to
\[
x_1(p',m) - x_1(p',m^s) \le 0.
\]
Hence the income effect is nonpositive:
\[
\Delta x_1^n \le 0.
\]
From Theorem \ref{thm:ch2-sub-neg}, the substitution effect also satisfies
\[
\Delta x_1^s \le 0.
\]
Therefore, by the Slutsky identity,
\[
\Delta x_1 = \Delta x_1^s + \Delta x_1^n \le 0.
\]
\end{proof}

\begin{examplebox}
\textbf{Computing the Slutsky effects explicitly.}

Use the demand system
\[
x_1(p_1,p_2,m)=\frac{m p_2}{p_1(p_1+p_2)}.
\]
Take
\[
p=(1,1),\qquad m=2,\qquad p'=(2,1).
\]
Original demand:
\[
x(p,m)=\left(1,1\right).
\]
Compensated expenditure:
\[
m^s = p' \cdot x(p,m)=2(1)+1(1)=3.
\]
Compensated demand for good 1:
\[
x_1(p',m^s)=\frac{3\cdot 1}{2(2+1)}=\frac12.
\]
Uncompensated new demand for good 1:
\[
x_1(p',m)=\frac{2\cdot 1}{2(2+1)}=\frac13.
\]

So:
\[
\Delta x_1^s = \frac12 - 1 = -\frac12,
\]
\[
\Delta x_1^n = \frac13 - \frac12 = -\frac16,
\]
and hence
\[
\Delta x_1
=
\Delta x_1^s + \Delta x_1^n
=
-\frac12 - \frac16
=
-\frac23.
\]

Interpretation:
\begin{itemize}
    \item the substitution effect is negative, as theory predicts;
    \item the income effect is also negative because good 1 is normal;
    \item therefore the total demand response is negative.
\end{itemize}
\end{examplebox}

\begin{examtip}
For price-change questions, keep the order of objects clear:
\begin{enumerate}
    \item old choice $x(p,m)$;
    \item new prices $p'$;
    \item compensated expenditure $m^s = p' \cdot x(p,m)$;
    \item compensated demand $x(p',m^s)$;
    \item uncompensated new demand $x(p',m)$.
\end{enumerate}
Most sign errors come from mixing up $m$ and $m^s$.
\end{examtip}

\subsection{Chapter Summary}

\begin{keyformula}[title={Key Formula: Lecture 2 Summary}]
\begin{align*}
x R^* x' &\iff \text{$x$ is revealed preferred to $x'$ through a chain of direct revealed preferences}, \\
x P^* x' &\iff \text{the chain contains at least one strict link}, \\
\text{GARP:}\quad x R^* x' &\implies \text{not }(x' P x), \\
m^s &= p' \cdot x(p,m), \\
\Delta x_1 &= x_1(p',m)-x_1(p,m), \\
\Delta x_1^s &= x_1(p',m^s)-x_1(p,m), \\
\Delta x_1^n &= x_1(p',m)-x_1(p',m^s), \\
\Delta x_1 &= \Delta x_1^s + \Delta x_1^n.
\end{align*}
If good 1 is normal and its price rises, then both $\Delta x_1^s \le 0$ and $\Delta x_1^n \le 0$, hence
\[
\Delta x_1 \le 0.
\]
\end{keyformula}

\begin{examtip}
The highest-yield ideas from Lecture 2 are:
\begin{itemize}
    \item GARP is the multi-observation revealed-preference test and rationalizability condition.
    \item Marshallian demand describes the \emph{total} response to price and expenditure changes.
    \item Gross substitutes/complements refer to cross-price effects; ordinary/Giffen refer to own-price effects.
    \item Normal/inferior refer to expenditure effects.
    \item The Slutsky identity decomposes a price effect into substitution and income effects.
    \item For normal goods, the law of demand follows immediately because both components are nonpositive.
\end{itemize}
\end{examtip}